\Needspace{18\baselineskip}
\section{Hole inventories and sharp shallow context closure}
\label{sec:context-closure}
All hole-label comparisons in this section compare complete typed labels: the
name and both interface lists. The occurrence quantifiers include every box depth.

\Needspace{24\baselineskip}
\subsection{Hole inventories, absence, and isomorphism}
\begin{definition}[Presence, absence, and disjointness]
Write $\mathsf{HasHole}(G,h)$ when some occurrence of $G$ looks up to $h$.
Write $\mathsf{HoleFree}(G)$ when no label has such an occurrence. Two graphs
have disjoint holes when no complete hole label occurs in both. The graphs may
have different external interfaces and depth bounds.
\end{definition}
\lean{Graph/ContextPlugging}{11}{12}{Linear.HasHole}
\lean{Graph/ContextPlugging}{14}{15}{Linear.HoleFree}
\lean{Graph/ContextPlugging}{18}{19}{Linear.HolesDisjoint}
Hole presence is equivalently witnessed by a fibre-coordinate occurrence with
the same observation. Transport a witness by the occurrence bijection for one
direction and by its inverse for the other, using observation preservation and
the inverse equation.
\lean{Graph/ContextPlugging}{21}{31}{Linear.hasHole_iff_occurrences}
A hole-free graph has unique holes vacuously: the first purported hole occurrence
already contradicts absence. Disjointness is symmetric by exchanging its two
presence premises. Any graph is disjoint from a hole-free graph because a hole
in the second graph would contradict its absence property.
\lean{Graph/ContextPlugging}{33}{36}{Linear.holeFree_unique}
\lean{Graph/ContextPlugging}{38}{40}{Linear.HolesDisjoint.symm}
\lean{Graph/ContextPlugging}{42}{44}{Linear.holeFree_disjoint}

Boxing neither creates nor removes a hole label from the inventory of the body.
A shallow box occurrence cannot be a hole. A deep occurrence is exactly an
address in its body with the same lookup, so it supplies a body witness. Conversely,
tag a body witness with the box's only outer edge. Consequently a hole-free body
gives a hole-free box: transfer any proposed hole witness to the body and contradict
its absence property.
\lean{Graph/ContextPlugging}{47}{58}{Linear.box_hasHole_iff}
\lean{Graph/ContextPlugging}{60}{63}{Linear.box_holeFree}

A label-preserving edge map preserves hole presence: map the fibre-coordinate
witness using the occurrence map and its observation equation. An injective such
edge map reflects unique holes: map two proposed equal-label occurrences into
the target, apply target uniqueness, and use occurrence-map injectivity.
\lean{Graph/ContextPlugging}{66}{71}{Linear.hasHole_map}
\lean{Graph/ContextPlugging}{74}{82}{Linear.uniqueHoles_of_embedding}
It follows that graph isomorphisms preserve and reflect hole inventories, and
preserve and reflect the context condition. Apply presence preservation to the
forward and inverse edge bijections for the first equivalence. Apply reflection
of uniqueness to the inverse bijection for forward preservation and to the
forward bijection for reflection. Label preservation is supplied by the
isomorphism equations with the appropriate equality direction.
\lean{Graph/ContextPlugging}{84}{87}{Linear.iso_hasHole_iff}
\lean{Graph/ContextPlugging}{89}{94}{Linear.iso_uniqueHoles_iff}
These results use the current isomorphism notion that fixes edge labels. They do
not yet extend it to isomorphic, independently renamed box contents.

\Needspace{24\baselineskip}
\subsection{Provenance and the abstract collision argument}
\begin{definition}[Provenance certificate]
For result $R$, host $C$, and replacement $H$, a provenance certificate is an
injective function
\[
 p:\operatorname{Occurrences}(R)\longrightarrow
 \operatorname{Occurrences}(C)+\operatorname{Occurrences}(H)
\]
that preserves hole observation in the chosen summand. Its left values identify
retained host occurrences and its right values identify inserted occurrences.
Different result occurrences must have different tagged origins.
\end{definition}
\lean{Graph/ContextPlugging}{98}{106}{Linear.Provenance}
This is a reusable intermediate interface. Unlike the earlier specimen laws,
the provenance certificates required for actual plugging are constructed below
from the concrete operations, rather than left as unproved assumptions.

Every result hole label belongs to the host or replacement inventory. Choose its
fibre-coordinate witness, inspect the origin tag, and use preservation to obtain
a witness in the indicated graph.
\lean{Graph/ContextPlugging}{113}{119}{Linear.Provenance.holes_subset}
A provenance certificate is \emph{compatible} when no inserted hole label is
observed at any retained host origin. Formally, if $p(a)=\mathsf{left}(c)$ and
$c$ observes $h$, then $H$ does not contain $h$. Host occurrences outside the
image of the left branch are irrelevant to this condition.
\lean{Graph/ContextPlugging}{122}{124}{Linear.Provenance.Compatible}

\begin{theorem}[Unique holes from compatible provenance]
If $C,H$ have unique holes and a provenance certificate for $R,C,H$ is compatible,
then $R$ has unique holes.
\end{theorem}
\begin{proof}
Take two result occurrences observing the same $h$ and inspect their origins.
If both origins are retained, host uniqueness identifies their host addresses.
If both are inserted, replacement uniqueness identifies their replacement
addresses. In either case the tagged origins agree and provenance injectivity
identifies the result occurrences. If one origin is retained and one inserted,
observation preservation gives a retained host observation of $h$ and a hole
witness for $h$ in $H$, contradicting compatibility. This covers all four tag
combinations and proves fibre-coordinate uniqueness, hence the original condition.
\end{proof}
\lean{Graph/ContextPlugging}{126}{151}{Linear.Provenance.unique_of_compatible}
Full disjointness of the host and replacement inventories implies compatibility:
any retained host occurrence is in particular a host occurrence. Substituting
this fact into the preceding theorem gives unique holes under full disjointness.
\lean{Graph/ContextPlugging}{154}{157}{Linear.Provenance.unique}

\Needspace{24\baselineskip}
\subsection{Concrete shallow provenance and a first context operation}
The provenance map for $C[e:=H]$ reads its outer edge tag. A surviving host edge
with fibre address $x$ maps to the corresponding old host edge with $x$;
an inserted edge maps to its replacement edge with $x$. In either case the
edge label, and hence the fibre, is definitionally unchanged.
\begin{proof}[Injectivity and observation preservation]
Mixed edge tags give different origin tags and cannot agree. Two retained
origins first identify the old edge values. Their surviving-edge subtype values
then agree because their inequality certificates are propositions; equality of
the dependent second components identifies their fibre addresses. Two inserted
origins are already the same replacement pair. In every case fibre observation
is unchanged by definition. Thus these data construct a provenance certificate
for any shallow replacement, including when the removed edge is a box.
\end{proof}
\lean{Graph/ContextPlugging}{163}{190}{Linear.shallowProvenance}
The general provenance theorem now proves context closure for shallow plugging
under full hole disjointness. A context-valued operation packages the filled
graph with this proof, taking certified host and replacement contexts and a
proof of disjointness. All graph incidence invariants still come from the earlier
certified net plugging construction.
\lean{Graph/ContextPlugging}{193}{198}{Linear.plug_uniqueHoles}
\lean{Graph/ContextPlugging}{200}{204}{Linear.Context.plug}

\Needspace{24\baselineskip}
\subsection{Exact inventory and the sharp freshness condition}
\begin{definition}[Retained hole label]
A retained hole label for removal of edge $e$ in $C$ is a hole observation in the
fibre of an outer edge $f\ne e$. This includes deep occurrences below $f$.
It excludes both the removed edge and every occurrence inside its label.
\end{definition}
\lean{Graph/ContextPlugging}{207}{209}{Linear.RetainedHole}
\begin{theorem}[Exact shallow hole inventory]
For each label $h$,
\[
 \mathsf{HasHole}(C[e:=H],h)
 \quad\Longleftrightarrow\quad
 \mathsf{RetainedHole}(C,e,h)\ \lor\ \mathsf{HasHole}(H,h).
\]
\end{theorem}
\begin{proof}
Express presence using fibre coordinates. A result occurrence has an outer edge
tag that is either a surviving host edge or an inserted edge. The first case
supplies exactly the retained witness; the second supplies a replacement witness.
Conversely, a retained witness gives a result occurrence at the left edge tag,
using its inequality proof to construct the surviving-edge subtype. A replacement
witness gives a result occurrence at the right tag. Their hole observations
are unchanged by definition.
\end{proof}
\lean{Graph/ContextPlugging}{212}{229}{Linear.plug_hasHole_iff}

Define sharp plugging freshness to mean that no retained host hole label occurs
in the replacement. This is weaker than requiring the replacement to be disjoint
from the entire original host: labels confined to the removed edge's region may
be reused.
\lean{Graph/ContextPlugging}{232}{235}{Linear.PlugFresh}
Sharp freshness implies uniqueness of the filled graph when both inputs have
unique holes. Apply the compatible-provenance theorem to shallow provenance.
Any result occurrence with a retained origin must have a surviving host edge
tag; its fibre is the retained witness required by sharp freshness. An inserted
edge cannot have a retained origin, so the other case is impossible.
\lean{Graph/ContextPlugging}{237}{249}{Linear.plug_uniqueHoles_of_fresh}

\begin{theorem}[Sharp freshness is necessary]
If $C[e:=H]$ has unique holes, sharp plugging freshness holds.
\end{theorem}
\begin{proof}
Suppose a retained host fibre observes $h$ and $H$ contains $h$. Form the two
result occurrences from their witnesses as in the inventory theorem. Result
uniqueness identifies them. Equality of their first components would identify
a surviving host edge tag with an inserted edge tag, which is impossible.
\end{proof}
\lean{Graph/ContextPlugging}{252}{261}{Linear.plug_fresh_necessary}
Consequently, for two contexts the filled graph has unique holes if and only if
sharp plugging freshness holds. Necessity is the preceding theorem, which did
not require uniqueness of either input; sufficiency uses the two input certificates.
\lean{Graph/ContextPlugging}{264}{269}{Linear.plug_uniqueHoles_iff_fresh}

\Needspace{24\baselineskip}
\subsection{Consumption of a selected hole and certified composition}
In a context, if the removed edge's shallow lookup is $h$, that label has no
retained occurrence elsewhere. Otherwise the retained occurrence and the selected
edge's shallow occurrence would be two occurrences of $h$. Context uniqueness
identifies their pairs, so their outer edges agree, contradicting the retained
edge's inequality certificate.
\lean{Graph/ContextPlugging}{272}{277}{Linear.removed_hole_not_retained}
It follows that the consumed label $h$ occurs in $C[e:=H]$ exactly when it occurs
in $H$. Use the exact inventory theorem and exclude its retained alternative
by the preceding result. Thus consumption removes the old occurrence; it does
not forbid the replacement from deliberately reintroducing that label.
\lean{Graph/ContextPlugging}{280}{287}{Linear.plug_consumed_hole_iff}
The operation $\mathsf{Context.plugFresh}$ packages shallow composition with this
sharp freshness premise. Its result graph is ordinary net plugging, and its
context certificate is the sharp sufficiency proof. It supports arbitrary finite
depth and arbitrary selected outer edges, rather than requiring that the selected
edge's label itself be a hole.
\lean{Graph/ContextPlugging}{289}{293}{Linear.Context.plugFresh}
